% Part: conditional-logics
% Chapter: minimal-change-semantics
% Section: sphere-models

\documentclass[../../../include/open-logic-section]{subfiles}

\begin{document}

\olfileid{con}{min}{sph}

\olsection{గోళ నమూనాలు}

ప్రతివాస్తవాలకు నియత అర్థవిచారం
ఇవ్వడానికి ఒక మార్గం లూయిస్
అనియత వివరణను గణిత నిర్మాణంగా
మార్చడం. అప్పుడు ఒక లోకం~$w$
చుట్టూ ఉన్న గోళాలు లోకాల
సమితులవుతాయి. గోళాలు ఒకదానిలో
మరొకటి ఉంటాయి కాబట్టి,~$w$
చుట్టూ ఉన్న లోకాల సమితులు
ఉపసమితి సంబంధం ప్రకారం రేఖీయ
క్రమంలో ఉండాలి.

\begin{defn}
  \emph{గోళ నమూనా} అనేది
  త్రయం~$\mModel{M} = \tuple{W, O, V}$;
  ఇక్కడ $W$ శూన్యం కాని లోకాల
  సమితి; $V\colon \PVar \to \Pow{W}$
  ఒక సత్యమూల్య నియామకం; $O\colon W
  \to \Pow{\Pow{W}}$ ప్రతి లోకం~$w$కు
  ఒక \emph{గోళ వ్యవస్థ}~$O_w$ను
  కేటాయిస్తుంది. ప్రతి $w$కు,
  $O_w$ లోకాల సమితుల సమితి;
  దానికి కింది షరతులు వర్తించాలి:
  \begin{enumerate}
  \item $O_w$,~$w$ వద్ద
    \emph{కేంద్రితమైనది}: $\{w\} \in O_w$.
  \item $O_w$ గోళాలు \emph{అంతర్నిహితమైనవి}:
    $S_1$, $S_2 \in O_w$
    అయినప్పుడల్లా $S_1
    \subseteq S_2$ లేదా $S_2
    \subseteq S_1$; అంటే $O_w$,
    ~$\subseteq$ ప్రకారం రేఖీయ
    క్రమంలో ఉంటుంది.
  \item $O_w$లోని గోళాల
    శూన్యం కాని ఉపకుటుంబాల
    సమ్మేళనాల కింద అది సంవృతం.
  \item $O_w$లోని గోళాల
    శూన్యం కాని ఉపకుటుంబాల
    ఛేదనాల కింద అది సంవృతం.
  \end{enumerate}
\end{defn}

$O_w$ వెనుక సహజ భావం: $w$
``చుట్టూ'' ఉన్న లోకాలు~$w$కు
ఎంత దూరంలో ఉన్నాయన్నదాని ప్రకారం
పొరలుగా అమరుతాయి. అత్యంత లోపలి
గోళం $w$ ఒక్కటే, అంటే
సమితి~$\{w\}$: మరే ఇతర
గోళంలోని లోకాలకన్నా $w$
తనకు తానే, అంటే~$w$కు,
దగ్గర. $S \subsetneq S'$
అయితే $S' \setminus S$లోని
లోకాలు~$S$లోని లోకాలకన్నా~$w$కు
దూరంగా ఉంటాయి: $S' \setminus S$
అనేది~$S$కు,~$S'$ బయట ఉన్న
లోకాలకు మధ్య ``పొర''.
ప్రత్యేకించి, గోళాలు వాటి బయటి
ఉపరితలం లోపలున్న అన్ని లోకాలనూ
కలిగి ఉన్నాయని భావించాలి;
అవి విడివిడి పొరలు మాత్రమే కావు.

\begin{figure}
\begin{center}
\begin{tikzpicture}[layerwidth=1.5,scale=.6]\tiny
  \spheresystem{3}
  \propositionintersect{0}{3}{60}{6}
  \path (0,0) node[world] {$w$};
  \spherepos{-30}{2}{node[world] {$w_2$}}
  \spherepos{220}{2}{node[world] {$w_3$}}
  \spherepos{90}{2}{node[world] {$w_1$}}
  \spherepos[shift={(.1,0)}]{10}{3}{node[world] {$w_5$}}
  \spherepos[shift={(.1,0)}]{-10}{3}{node[world] {$w_6$}}
  \spherepos{225}{3}{node[world] {$w_4$}}
  \spherepos{20}{4}{node[world] {$w_7$}}
  \path (5,0) node[left] {$p$};
\end{tikzpicture}
\caption{గోళ నమూనా చిత్రం}
\ollabel{fig:sphere-model}
\end{center}
\end{figure}

\olref{fig:sphere-model}లోని చిత్రం
$W = \{w, w_1, \dots, w_7\}$,
$V(p) = \{w_5, w_6,
w_7\}$ గల గోళ నమూనాను
సూచిస్తుంది. అత్యంత లోపలి
గోళం $S_1 = \{w\}$.
$w$కు అత్యంత సమీప లోకాలు
$w_1, w_2, w_3$ కాబట్టి,
తదుపరి పెద్ద గోళం
$S_2 = \{w, w_1,
w_2, w_3\}$. ఇంకా బయట
ఉన్న లోకాలు $w_4$, $w_5$,
$w_6$; అందువల్ల అత్యంత
బయటి గోళం $S_3 = \{w,
w_1, \dots, w_6\}$.
$w$ చుట్టూ గోళ వ్యవస్థ
$O_w = \{S_1, S_2, S_3\}$.
లోకం~$w_7$,~$w$ చుట్టూ
ఉన్న ఏ గోళంలోనూ లేదు.
$p$~సత్యమయ్యే అత్యంత
సమీప లోకాలు $w_5$,~$w_6$;
కాబట్టి $p$ సత్యమయ్యే
లోకాన్ని కలిగిన అత్యంత చిన్న
గోళం~$S_3$.

గోళ నమూనా~$\mModel M$లో
లోకం~$w$ వద్ద సూత్రం~$!A$
సంతృప్తిని, అంటే
$\mSat{M}{!A}[w]$ను
నిర్వచించడానికి, మోడల్
\tetoken{సూత్రాల}{!!{formula}s}
నిర్వచనానికి $!B \cif !C$
కోసం ఒక షరతును చేరుస్తాం:

\begin{defn}
  $\mSat{M}{!B \cif !C}[w]$
  అప్పుడే కింది రెండింటిలో
  ఏదో ఒకటి జరుగుతుంది:
  \begin{enumerate}
  \item\ollabel{sphere-vac} అన్ని $u \in \bigcup O_w$లకు
    $\mSat/{M}{!B}[u]$, లేదా
  \item\ollabel{sphere-nonvac} ఏదో $S \in O_w$కు,
    \begin{enumerate}
      \item ఏదో $u \in
        S$కు $\mSat{M}{!B}[u]$,
        మరియు
      \item అన్ని $v \in S$లకు
        $\mSat/{M}{!B}[v]$ లేదా
        $\mSat{M}{!C}[v]$.
    \end{enumerate}
  \end{enumerate}
\end{defn}

ఈ నిర్వచనం ప్రకారం,
$\mSat{M}{!B \cif !C}[w]$
అప్పుడే: పూర్వపక్షం~$!B$,
~$w$ చుట్టూ ఉన్న గోళాలన్నిటిలో
ప్రతిచోట అసత్యం; లేదా $!B$
సత్యమయ్యే గోళం~$S$ ఏదో ఉంది,
ఆ ``$!B$ను కలిగిన'' గోళంలోని
అన్ని లోకాల వద్ద భౌతిక సోపాధికం
$!B \lif !C$ సత్యం. సహజ
వివరణను బట్టి అనుకున్నట్లు,
నిర్వచనంలో $S$ తప్పనిసరిగా
$!B$ను కలిగిన \emph{అత్యంత లోపలి}
గోళం కావాలని కోరలేదని గమనించండి.
కానీ~\olref{sphere-nonvac}లోని
షరతు ఏదో గోళం~$S$కు
వర్తిస్తే,~$S$ లోపల ఉండి
పూర్వపక్షం సత్యమయ్యే ప్రతి గోళానికీ
వర్తిస్తుంది; అందువల్ల అలాంటి
అత్యంత లోపలి గోళం ఉన్నప్పుడు
దానికీ వర్తిస్తుంది.
\sourcecorrection{OLTECNTSPH-001}{మూల వివరణలో పెద్ద గోళంలోని షరతు అది కలిగిన అన్ని చిన్న గోళాలకూ వర్తిస్తుందని చెప్పింది. చిన్న గోళంలో పూర్వపక్షం సత్యమయ్యే లోకం లేకపోవచ్చు; కాబట్టి పూర్వపక్షాన్ని కలిగిన చిన్న గోళాలకే వాదనను పరిమితం చేశాం, అత్యంత లోపలి గోళం ఉంటే అనే షరతునూ స్పష్టం చేశాం.}

$!B$ను కలిగిన అత్యంత
లోపలి గోళం \emph{ఉండాలని}
గోళ నమూనాల నిర్వచనం
కోరదని కూడా గమనించండి:
$!B$ను కలిగిన గోళాల
అనంత క్రమం $S_1
\supsetneq S_2 \supsetneq
\dots \supsetneq
\{w\}$ ఉండవచ్చు;
అప్పుడు $!B$ను కలిగిన
అత్యంత లోపలి గోళం
ఉండదు. అలాంటి సందర్భంలో,
ఏదో~$i$కు $S_i$,
$S_{i+1}$, \dots గోళాలన్నిటిలో
$!B \lif !C$ సత్యమైనప్పుడే
$\mSat{M}{!B \cif !C}[w]$.

\end{document}
